Наконец, мы готовы сформулировать и доказать вторую версию усиленного закона больших чисел. Напомним сначала определение понятия одинаковой распределенности.

\mkdef{Одинаковая распределенность}{id_dist}{
    Случайные величины $\{\xi_\alpha, \alpha \in \mathfrak{A}\}$ называются одинаково распределенными, если все они имеют одинаковую функцию распределения. Обозначение: $\xi \stackrel{d}{=} \eta$ ($\xi$ и $\eta$ одинаково распределены).
}

Из формул подсчета математических ожиданий мы помним, что если $\xi \stackrel{d}{=} \eta$, то $\E g(\xi) = \E g(\eta)$ для любой борелевской функции $g(x)$.

\mkthm{Усиленный закон больших чисел в форме Этемади}{slln_etemadi}{
    Пусть $\{\xi_n, n \in \mathbb{N}\}$ — последовательность попарно независимых одинаково распределенных случайных величин, причем $\E |\xi_1| < \infty$. Обозначим $S_n = \xi_1 + \dots + \xi_n$. Тогда
    $$ \frac{S_n}{n} \toAS \E \xi_1 \quad \text{при } n \to \infty. $$
}

\mkproof{
    Рассмотрим разложения $\xi_n = \xi_n^+ - \xi_n^-$. Тогда последовательности $\{\xi_n^+, n \in \mathbb{N}\}$ и $\{\xi_n^-, n \in \mathbb{N}\}$ также удовлетворяют условию теоремы, то есть это последовательности попарно независимых одинаково распределенных случайных величин с конечным математическим ожиданием. Если мы докажем утверждение для них, то оно будет выполнено и для последовательности $\{\xi_n, n \in \mathbb{N}\}$. Тем самым, достаточно рассматривать только случай неотрицательных случайных величин. Считаем далее, что $\xi_n \ge 0$ для всех $n$.

    По условию $\E \xi_1 < \infty$, поэтому из утверждения об эквивалентности сходимости ряда и конечности матожидания, а также одинаковой распределенности с.в. получаем, что
    $$ \sum_{n=1}^\infty \P(\xi_n \ge n) = \sum_{n=1}^\infty \P(\xi_1 \ge n) < \infty. $$

    Тогда согласно лемме Бореля–Кантелли выполнено
    $$ \P(\{\xi_n \ge n \text{ б.ч.}\}) = 0. $$

    Следовательно, с вероятностью 1 для всех $n$, кроме лишь конечного числа, выполнено $\{\xi_n < n\}$. Обозначим это событие через $\Omega'$.

    Обозначим $\tilde{\xi}_n = \xi_n \I\{\xi_n < n\}$. Это будут по-прежнему независимые, но вообще говоря не одинаково распределенные случайные величины. Заметим, что для любого фиксированного $\omega \in \Omega'$ выполнено $\tilde{\xi}_n(\omega) = \xi_n(\omega)$ для всех $n$, кроме лишь конечного числа значений. Значит, для любого $\omega \in \Omega'$
    $$ \frac{\xi_1(\omega) + \dots + \xi_n(\omega)}{n} \to \E \xi_1 \iff \frac{\tilde{\xi}_1(\omega) + \dots + \tilde{\xi}_n(\omega)}{n} \to \E \xi_1. $$

    Поймем теперь, как ведет себя последовательность $\E \tilde{\xi}_n$. Пользуясь одинаковой распределенностью и теоремой Лебега о мажорируемой сходимости, получаем:
    $$ \E \tilde{\xi}_n = \E(\xi_n \I\{\xi_n < n\}) = \E(\xi_1 \I\{\xi_1 < n\}) \to \E \xi_1 \quad \text{при } n \to \infty. $$

    Отсюда по лемме Тёплица
    $$ \frac{1}{n} \sum_{k=1}^n \E \tilde{\xi}_k \to \E \xi_1, $$
    значит,
    $$ \frac{\tilde{\xi}_1 + \dots + \tilde{\xi}_n}{n} \to \E \xi_1 \iff \frac{(\tilde{\xi}_1 - \E \tilde{\xi}_1) + \dots + (\tilde{\xi}_n - \E \tilde{\xi}_n)}{n} \to 0. $$

    Покажем, что правая часть выполнена с вероятностью 1 для некоторой специальной подпоследовательности. Обозначим $\tilde{S}_n = \tilde{\xi}_1 + \dots + \tilde{\xi}_n$. Зафиксируем произвольное $\alpha > 1$ и рассмотрим последовательность $\{k_m, m \in \mathbb{N}\}$, где $k_m = [\alpha^m]$. Проверим, что $(\tilde{S}_{k_m} - \E \tilde{S}_{k_m}) / k_m \toAS 0$ при $m \to \infty$. Для произвольного $\varepsilon > 0$ имеем:

    \begin{align*}
        \sum_{m=1}^\infty \P \left( \left| \frac{\tilde{S}_{k_m} - \E \tilde{S}_{k_m}}{k_m} \right| \ge \varepsilon \right) &\le | \text{неравенство Чебышева} | \le \sum_{m=1}^\infty \frac{\D \tilde{S}_{k_m}}{\varepsilon^2 k_m^2} = | \text{нез. } \tilde{\xi}_n | = \\
        &= \sum_{m=1}^\infty \frac{1}{\varepsilon^2 k_m^2} \sum_{j=1}^{k_m} \D \tilde{\xi}_j \le \sum_{m=1}^\infty \frac{1}{\varepsilon^2 k_m^2} \sum_{j=1}^{k_m} \E(\tilde{\xi}_j)^2 = \\
        &= \frac{1}{\varepsilon^2} \sum_{j=1}^\infty \E(\tilde{\xi}_j)^2 \sum_{m \colon k_m \ge j} \frac{1}{k_m^2}.
    \end{align*}

    Далее, воспользуемся тем, что последовательность $k_m$ ведет себя фактически как геометрическая прогрессия. Тогда сумма $\sum_{m \colon k_m \ge j} \frac{1}{k_m^2}$ ведет себя по порядку как свое первое слагаемое, которое заведомо не превосходит $1/j^2$. Следовательно, существует такая величина $C = C(\alpha) > 0$, что

    \begin{align*}
        \sum_{m=1}^\infty \P \left( \left| \frac{\tilde{S}_{k_m} - \E \tilde{S}_{k_m}}{k_m} \right| \ge \varepsilon \right) &\le \frac{C}{\varepsilon^2} \sum_{j=1}^\infty \frac{\E(\tilde{\xi}_j)^2}{j^2} = \\
        &= \frac{C}{\varepsilon^2} \sum_{j=1}^\infty \frac{1}{j^2} \E \left( \xi_j^2 \I\{\xi_j < j\} \right) = | \text{одинак. распред.} | = \\
        &= \frac{C}{\varepsilon^2} \sum_{j=1}^\infty \frac{1}{j^2} \E \left( \xi_1^2 \I\{\xi_1 < j\} \right) = \\
        &= \frac{C}{\varepsilon^2} \sum_{j=1}^\infty \frac{1}{j^2} \sum_{k=1}^j \E \left( \xi_1^2 \I\{k - 1 \le \xi_1 < k\} \right) = \\
        &= \frac{C}{\varepsilon^2} \sum_{k=1}^\infty \E \left( \xi_1^2 \I\{k - 1 \le \xi_1 < k\} \right) \sum_{j=k}^\infty \frac{1}{j^2} \le \Big| \text{оценим } \sum_{j=k}^\infty \frac{1}{j^2} \le \frac{2}{k} \Big| \le \\
        &\le \frac{2C}{\varepsilon^2} \sum_{k=1}^\infty \frac{1}{k} \E \left( \xi_1^2 \I\{k - 1 \le \xi_1 < k\} \right) \le \\
        &\le \frac{2C}{\varepsilon^2} \sum_{k=1}^\infty \E \left( \xi_1 \I\{k - 1 \le \xi_1 < k\} \right) = \\
        &= | \text{монотонная сходимость} | = \frac{2C}{\varepsilon^2} \E \xi_1 < \infty.
    \end{align*}

    Достаточное условие сходимости почти наверное выполнено и мы доказали, что $(\tilde{S}_{k_m} - \E \tilde{S}_{k_m}) / k_m \toAS 0$, а значит, $\tilde{S}_{k_m} / k_m \toAS \E \xi_1$.
    
    Осталось завершить доказательство. В силу неотрицательности случайных величин $\tilde{\xi}_n$ последовательность $\tilde{S}_n$ не убывает. Если $k_{m-1} < n \le k_m$, то будут выполнены неравенства
    $$ \frac{\tilde{S}_{k_{m-1}}}{k_{m-1}} \cdot \frac{k_{m-1}}{k_m} = \frac{\tilde{S}_{k_{m-1}}}{k_m} \le \frac{\tilde{S}_n}{n} \le \frac{\tilde{S}_{k_m}}{k_{m-1}} = \frac{\tilde{S}_{k_m}}{k_m} \cdot \frac{k_m}{k_{m-1}}. $$
    Применяя доказанную сходимость и учитывая, что $\frac{k_m}{k_{m-1}} \to \alpha$ при $m \to \infty$, получаем, что с вероятностью 1
    $$ \frac{1}{\alpha} \E \xi_1 \le \varliminf_{n} \frac{\tilde{S}_n}{n} \le \varlimsup_{n} \frac{\tilde{S}_n}{n} \le \alpha \E \xi_1. $$
    В силу произвольности $\alpha > 1$ это означает, что $\tilde{S}_n / n \toAS \E \xi_1$. Осталось вспомнить, что пределы $\tilde{S}_n / n$ и $S_n / n$ совпадают с вероятностью 1.
}

\vspace{1.5em}
В качестве мгновенного следствия из теоремы  мы получаем знаменитый усиленный закон больших чисел в форме Колмогорова.

\mkcor{Усиленный закон больших чисел в форме Колмогорова}{slln_kolm}{
    Пусть $\{\xi_n, n \in \mathbb{N}\}$ — последовательность независимых одинаково распределенных случайных величин (н.о.р.с.в.), причем $\E |\xi_1| < \infty$. Тогда
    $$ \frac{\xi_1 + \dots + \xi_n}{n} \toAS \E \xi_1 \quad \text{при } n \to \infty. $$
}

\vspace{1.5em}
Можно пойти дальше и задаться вопросом о необходимости конечности математического ожидания для сходимости сумм независимых одинаково распределенных случайных величин. Несложно понять, однако, что условие наличия конечного математического ожидания действительно необходимо.

\mkcor{Критерий выполнимости УЗБЧ для н.о.р.с.в.}{crit_slln}{
    Пусть $\{\xi_n, n \in \mathbb{N}\}$ — последовательность независимых одинаково распределенных случайных величин (н.о.р.с.в.). Обозначим $S_n = \xi_1 + \dots + \xi_n$. Тогда последовательность $S_n / n$ сходится почти наверное (к произвольной случайной величине) тогда и только тогда, когда $\E \xi_1$ конечно.
}

\mkproof{
    ($\Rightarrow$) Пусть $S_n/n \toAS \eta$. Тогда
    $$ \frac{\xi_n}{n} = \frac{S_n}{n} - \frac{S_{n-1}}{n} = \frac{S_n}{n} - \frac{S_{n-1}}{n-1} \cdot \frac{n}{n-1} \toAS 0. $$
    Следовательно, сходится ряд из вероятностей $\P(|\xi_n| \ge n)$. Действительно, если ряд расходится, то по лемме Бореля–Кантелли мы получаем, что с вероятностью 1 будет выполнено бесконечное число событий $\{|\xi_n| \ge n\}$, что противоречит сходимости к нулю последовательности $\xi_n/n$. Но в силу одинаковой распределенности
    $$ \sum_{n=1}^\infty \P(|\xi_n| \ge n) = \sum_{n=1}^\infty \P(|\xi_1| \ge n) < \infty. $$
    Из критерия конечности матожидания получаем, что $\E |\xi_1|$ конечно.

    ($\Leftarrow$) \hyperref[cor:slln_kolm]{Следствие 4.4} показывает, что $\eta = \E \xi_1$.
}

